\section{A numerical CM bound for every split class-number-one datum}

\begin{theorem}[HD-CM6: an explicit split-sixfold bound]
\label{HD-CM6}
Let $K=\mathbf Q(\sqrt{-m})$ have class number one, and define the
integers $L,d=4mL^2$ as in SF-QC01. Let $(A,\theta)$ be a split
Weil sixfold with integral compatible maximal-order action, and
$\Delta=\int_A\theta^6/6!$. Then, using the supplied Perry
deformation input with its stated preprint status,
\[
 \boxed{\delta_W(A,\theta)\le (2d\Delta)^6D_{\rm CM}(m)
                         \le2^{324}d^{57}\Delta^6.}
\]
The same bound holds uniformly on the entire fixed polarized family.
If only $\mathbf Z+f_t\mathcal O_K$ acts integrally, replace
$\Delta$ by $f_t^{12}\Delta$, giving the safe last bound
$2^{324}d^{57}f_t^{72}\Delta^6$.
\end{theorem}
\begin{proof}
We make the source-to-hyperbolic lattice map explicit; a varying
isogeny degree is not being used as an uncontrolled parameter.
The source of SF-QC05 is $J\times\widehat J$, identified with
$J^2$ by its principal polarization, with
$h=3(d\Theta_x+\Theta_y)$ and
\[
 u=\begin{pmatrix}0&-1\\d&0\end{pmatrix},\qquad u^2=-d.
\]
Choose an integral symplectic basis $a_1,a_2,a_3,b_1,b_2,b_3$
of $H_1(J,\mathbf Z)$, and use its vectors in the first factor
as a $K$-basis of the full rational homology. The second-factor
vector corresponding to $v$ is $(u/d)v$. Thus the integral source
lattice in this $K$-basis is
\[
 \Lambda_s=(\mathbf Z+(u/d)\mathbf Z)^6.
\]
Write $\mathcal O_K=\mathbf Z[\tau]$, where
$\tau=a+u/f_s$. If $m\not\equiv3\pmod4$, take
$a=0,f_s=2L$; if $m\equiv3\pmod4$, take
$a=1/2,f_s=4L$. In both cases $f_s\mid d$.
For the normalized Hermitian form
$h_K(v,w)=E_h(v,\tau w)-\bar\tau E_h(v,w)$,
the matrix in the displayed $K$-basis is
\[
 H_s=\begin{pmatrix}0&kI_3\\-kI_3&0\end{pmatrix},
 \qquad k=3du/f_s\in\mathcal O_K,
 \qquad\bar k=-k.
\]
Indeed $u$ takes a first-factor vector into the second factor,
which is $E_h$-orthogonal to the first, while the alternating form
on the first factor is $3dE_\Theta$. This proves the formula
including its normalization factor.

Let $U^3$ denote the standard hyperbolic Hermitian lattice on
$\mathcal O_K^6$. The matrix
$F=\operatorname{diag}(I_3,kI_3)$ satisfies
$F^*U^3F=H_s$. Moreover $dF\Lambda_s\subset\mathcal O_K^6$:
the first three coordinates lie in $d\mathbf Z+u\mathbf Z$,
and the last three in $dk\mathbf Z+ku\mathbf Z$.
Consequently, after transporting the complex structure, $dF$
defines an actual $K$-linear isogeny from the source to the
standard hyperbolic polarized abelian variety. Its pullback of
the target polarization is $d^2h$. A source threefold detector
of $h$-degree at most $D_{\rm CM}(m)$ therefore pushes to an
effective detector of degree at most $d^6D_{\rm CM}(m)$.
The rational isometry identifies the full period domains;
intersecting the arithmetic lattice stabilizers gives the finite
level correspondence. Thus this assertion holds throughout the
hyperbolic component, not just at its CM point.

The target lattice is free over $\mathcal O_K$ by the
class-number-one hypothesis. Intersect a maximal rational
isotropic subspace with it and extend a basis of that saturated
rank-three submodule to an integral basis. Its Hermitian matrix
is $\bigl(\begin{smallmatrix}0&B\\B^*&C\end{smallmatrix}\bigr)$,
with $\operatorname{Nm}(\det B)=\Delta$. The explicit substitution
in SF-04a changes this matrix to $U^3$ and has all denominators
dividing $2\Delta$ in both directions. That algebraic calculation
uses only a principal ideal domain and complex conjugation, so it
applies to the present ring, including the non-Euclidean
class-number-one rings. Multiplication by $2\Delta$ gives an
isogeny from the hyperbolic datum to the target with pullback
polarization $(2\Delta)^2\theta_U$. Pushforward therefore costs
at most $(2\Delta)^6$. Combining the two maps proves the first
bound, and SF-QC05 gives the second.

For a fixed smaller order, write $A=V/\Lambda$ and put
$\Lambda'=f_t\mathcal O_K\Lambda$. The inclusions
$f_t\Lambda\subset\Lambda'\subset\Lambda$ give an index
$I\le f_t^{12}$. The isogeny $V/\Lambda'\to A$ has pulled-back
polarization product $I\Delta$. Apply the maximal-order bound
and push forward; projection formula preserves the cycle degree
relative to these specified polarizations and preserves detection.
Monotonicity in $\Delta$ gives the conductor assertion.
\end{proof}

\begin{corollary}[HD-CM4: all fourfold determinants by compensation]
\label{HD-CM4}
For every class-number-one imaginary quadratic $K$, the same
explicit $d$ and $c_K=1+\operatorname{Nm}(\tau)$ give, on every
fourfold determinant and integral maximal-order polarization datum,
\[
 \delta_W(A,\theta)\le
 c_K(2d\Delta^2)^6D_{\rm CM}(m)
 \le c_K2^{324}d^{57}\Delta^{12}.
\]
For a fixed conductor $f_t$, the last expression is multiplied
by $f_t^{96}$. Smaller previously proved bounds are retained
when available. These conclusions use Perry as in SF-QC04;
no Markman, Severn or 032 theorem is a premise.
\end{corollary}
\begin{proof}
Apply HD-COMP with the explicit split-sixfold function HD-CM6.
The conductor statement uses the fourfold lattice index
$I\le f_t^8$ in the same pushforward argument.
\end{proof}
